\documentclass[10pt,a4paper]{amsart}
\usepackage{amsmath,amssymb,amsthm,graphicx}
% For commutative diagrams we use Xy-pic
\usepackage[all]{xy}

% We use 2cell for 2-commutative diagrams.
\xyoption{2cell}
% Theorem environments.
%
\theoremstyle{plain}
\newtheorem{theorem}[subsection]{Theorem}
\newtheorem{proposition}[subsection]{Proposition}
\newtheorem{lemma}[subsection]{Lemma}

\theoremstyle{definition}
\newtheorem{definition}[subsection]{Definition}
\newtheorem{example}[subsection]{Example}
\newtheorem{exercise}[subsection]{Exercise}
\newtheorem{situation}[subsection]{Situation}

\theoremstyle{remark}
\newtheorem{remark}[subsection]{Remark}
\newtheorem{remarks}[subsection]{Remarks}

\numberwithin{equation}{subsection}

% Macros
%
\def\lim{\mathop{\mathrm{lim}}\nolimits}
\def\colim{\mathop{\mathrm{colim}}\nolimits}
\def\Spec{\mathop{\mathrm{Spec}}}
\def\Hom{\mathop{\mathrm{Hom}}\nolimits}
\def\Ext{\mathop{\mathrm{Ext}}\nolimits}
\def\SheafHom{\mathop{\mathcal{H}\!\mathit{om}}\nolimits}
\def\SheafExt{\mathop{\mathcal{E}\!\mathit{xt}}\nolimits}
\def\Sch{\mathit{Sch}}
\def\Mor{\mathop{\mathrm{Mor}}\nolimits}
\def\Ob{\mathop{\mathrm{Ob}}\nolimits}
\def\Sh{\mathop{\mathit{Sh}}\nolimits}
\def\NL{\mathop{N\!L}\nolimits}
\def\CH{\mathop{\mathrm{CH}}\nolimits}
\def\proetale{{pro\text{-}\acute{e}tale}}
\def\etale{{\acute{e}tale}}
\def\QCoh{\mathit{QCoh}}
\def\Ker{\mathop{\mathrm{Ker}}}
\def\Im{\mathop{\mathrm{Im}}}
\def\Coker{\mathop{\mathrm{Coker}}}
\def\Coim{\mathop{\mathrm{Coim}}}

% Boxtimes
%
\DeclareMathSymbol{\boxtimes}{\mathbin}{AMSa}{"02}

%
% Macros for moduli stacks/spaces
%
\def\QCohstack{\mathcal{QC}\!\mathit{oh}}
\def\Cohstack{\mathcal{C}\!\mathit{oh}}
\def\Spacesstack{\mathcal{S}\!\mathit{paces}}
\def\Quotfunctor{\mathrm{Quot}}
\def\Hilbfunctor{\mathrm{Hilb}}
\def\Curvesstack{\mathcal{C}\!\mathit{urves}}
\def\Polarizedstack{\mathcal{P}\!\mathit{olarized}}
\def\Complexesstack{\mathcal{C}\!\mathit{omplexes}}
% \Pic is the operator that assigns to X its picard group, usage \Pic(X)
% \Picardstack_{X/B} denotes the Picard stack of X over B
% \Picardfunctor_{X/B} denotes the Picard functor of X over B
\def\Pic{\mathop{\mathrm{Pic}}\nolimits}
\def\Picardstack{\mathcal{P}\!\mathit{ic}}
\def\Picardfunctor{\mathrm{Pic}}
\def\Deformationcategory{\mathcal{D}\!\mathit{ef}}
\usepackage{cleveref}
\crefname{originalproposition}{original proposition}{original propositions}
\Crefname{originalproposition}{Original proposition}{Original propositions}
\makeatletter
\newlabel{proposition-etale-morphisms}{{03PC}{source}{}{sourcepointer-1}{}}
\newlabel{proposition-etale-morphisms@cref}{{[originalproposition][0][]03PC}{[1][0][]source}{}{}{}}
\makeatother
\crefname{originalremark}{original remark}{original remarks}
\Crefname{originalremark}{Original remark}{Original remarks}
\makeatletter
\newlabel{remark-etale-between-etale}{{03PP}{source}{}{sourcepointer-2}{}}
\newlabel{remark-etale-between-etale@cref}{{[originalremark][0][]03PP}{[1][0][]source}{}{}{}}
\makeatother
\title{The category of etale neighbourhoods of a geometric point is cofiltered}
\author{The Stacks Project Authors}
\date{stacks-project:etale-cohomology.tex:local-snapshot-20261011; GNU FDL 1.2 or later}
\begin{document}
\maketitle
\noindent\textit{Editorial note.} All delivered mathematical text is the author's own, copied from the pinned Stacks Project chapter file etale-cohomology.tex; no statement or proof step was rewritten, reordered or removed. The author's old-style \textbackslash{}not = and the arrow label \textasciicircum{}-\textbackslash{}Delta in the displayed square are reproduced verbatim.\par\medskip
\noindent\textit{Editorial note.} The preamble is the chapter's own native preamble excerpt used for this source: the Xy-pic setup at lines 14-18 and the theorem declarations with the mathematical macros at lines 154-217 of sources/stacks-etale-cohomology/preamble.tex. The quoted class line and the xr-hyper externaldocument block are not included because the assembler supplies a plain amsart wrapper; no macro was invented or redefined.\par\medskip
\noindent\textit{Editorial note.} Both referenced results live in the same chapter and are located at their \textbackslash{}label line inside the pinned chapter file.\par\medskip
\noindent\textit{Allowed standard prerequisites (editorial).}\begin{itemize}
\item Etale neighbourhood of a geometric point of a scheme. Same chapter, the definition introducing an etale neighbourhood (U, ubar) of a geometric point sbar of S, with U etale over S. Supplies the objects of the category whose cofilteredness is asserted; the two constructions are the counted content.
\item Permanence of etale morphisms under base change and composition, and between etale schemes. Same chapter, Proposition labelled proposition-etale-morphisms (source line 2592, Stacks tag 03PC) and Remark labelled remark-etale-between-etale (source line 3230, Stacks tag 03PP); invoked at proof lines 3279-3281, 3294-3296 and 3297-3298. Supply base change, composition and the criterion that a morphism between schemes etale over a base is etale; the fibre product constructions that use them are retained in the counted span.
\end{itemize}\medskip
\ifdefined\hypertarget\hypertarget{sourcepointer-1}{}\fi
\noindent\textit{Source reference.} Original proposition 03PC: sources/stacks-etale-cohomology/pinned-original-source.tex, source lines 2592--2592. Same-chapter pointer invoked at proof lines 3281 and 3296 for permanence of etale morphisms under base change and composition; the two fibre product constructions and the nonemptyness check are retained verbatim in the counted span.\par\medskip
\ifdefined\hypertarget\hypertarget{sourcepointer-2}{}\fi
\noindent\textit{Source reference.} Original remark 03PP: sources/stacks-etale-cohomology/pinned-original-source.tex, source lines 3230--3230. Same-chapter pointer invoked at proof line 3298 for the criterion that a morphism between schemes etale over S is etale, applied to (h\_1, h\_2); the base change and composition steps that follow are retained in the counted span.\par\medskip
\begin{lemma}
\label{lemma-cofinal-etale}
Let $S$ be a scheme, and let $\overline{s}$ be a geometric point of $S$.
The category of \'etale neighborhoods is cofiltered. More precisely:
\begin{enumerate}
\item Let $(U_i, \overline{u}_i)_{i = 1, 2}$ be two \'etale neighborhoods of
$\overline{s}$ in $S$. Then there exists a third \'etale neighborhood
$(U, \overline{u})$ and morphisms
$(U, \overline{u}) \to (U_i, \overline{u}_i)$, $i = 1, 2$.
\item Let $h_1, h_2: (U, \overline{u}) \to (U', \overline{u}')$ be two
morphisms between \'etale neighborhoods of $\overline{s}$. Then there exist an
\'etale neighborhood $(U'', \overline{u}'')$ and a morphism
$h : (U'', \overline{u}'') \to (U, \overline{u})$
which equalizes $h_1$ and $h_2$, i.e., such that
$h_1 \circ h = h_2 \circ h$.
\end{enumerate}
\end{lemma}

\begin{proof}
For part (1), consider the fibre product $U = U_1 \times_S U_2$.
It is \'etale over both $U_1$ and $U_2$ because \'etale morphisms are
preserved under base change, see
Proposition \ref{proposition-etale-morphisms}.
The map $\overline{s} \to U$ defined by $(\overline{u}_1, \overline{u}_2)$
gives it the structure of an \'etale neighborhood mapping to both
$U_1$ and $U_2$. For part (2), define $U''$ as the fibre product
$$
\xymatrix{
U'' \ar[r] \ar[d] & U \ar[d]^{(h_1, h_2)} \\
U' \ar[r]^-\Delta & U' \times_S U'.
}
$$
Since $\overline{u}$ and $\overline{u}'$ agree over $S$ with $\overline{s}$,
we see that $\overline{u}'' = (\overline{u}, \overline{u}')$ is a geometric
point of $U''$. In particular $U'' \not = \emptyset$.
Moreover, since $U'$ is \'etale over $S$, so is the fibre product
$U'\times_S U'$ (see
Proposition \ref{proposition-etale-morphisms}).
Hence the vertical arrow $(h_1, h_2)$ is \'etale by
Remark \ref{remark-etale-between-etale} above.
Therefore $U''$ is \'etale over $U'$ by base change, and hence also
\'etale over $S$ (because compositions of \'etale morphisms are \'etale).
Thus $(U'', \overline{u}'')$ is a solution to the problem.
\end{proof}

\end{document}
